\eventsection{CONSIDERA\c{C}\~OES FINAIS}

Este trabalho apresentou um sistema inteligente embarcado para monitoramento da qualidade da \'agua no contexto amaz\^onico, com aplica\c{c}\~ao no Lago do Serpa, em Itacoatiara, Amazonas. A proposta mostrou ser capaz de operar de forma offline, prever pH, oxig\^enio dissolvido e turbidez e executar processamento em diferentes plataformas de borda, incluindo ESP32, Raspberry Pi e smartphones Android. Os resultados permitem concluir que a solu\c{c}\~ao \'e operacionalmente vi\'avel e socialmente relevante para regi\~oes com baixa infraestrutura de comunica\c{c}\~ao. Seu uso potencial em comunidades ribeirinhas e em a\c{c}\~oes de acompanhamento ambiental pode contribuir para monitoramento cont\'inuo, apoio \`a tomada de decis\~ao e fortalecimento de estrat\'egias de preserva\c{c}\~ao e sa\'ude p\'ublica.

Tamb\'em se destaca que a proposta valoriza um caminho de inova\c{c}\~ao adequado \`a realidade regional. Em vez de depender de estruturas complexas e de alto custo, o sistema busca entregar intelig\^encia computacional em um formato compat\'ivel com o trabalho de campo. Essa escolha torna a solu\c{c}\~ao mais pr\'oxima de cen\'arios reais de uso e mais alinhada com demandas de territ\'orios remotos da Amaz\^onia e de regi\~oes marcadas por baixa conectividade.

A contribui\c{c}\~ao principal deste artigo est\'a justamente nesse deslocamento de foco: do detalhamento da modelagem para a demonstra\c{c}\~ao de aplicabilidade. Ao enfatizar o uso pr\'atico, a opera\c{c}\~ao offline e a relev\^ancia territorial da proposta, o trabalho se apresenta como uma contribui\c{c}\~ao voltada \`a ado\c{c}\~ao de tecnologia de baixo custo em monitoramento ambiental amaz\^onico. Os resultados comparativos tamb\'em mostram que a defesa da arquitetura espiculada deve ser feita com precis\~ao: ela n\~ao foi a melhor baseline em erro m\'edio nesta base, mas se mostrou competitiva e, ao mesmo tempo, compat\'ivel com execu\c{c}\~ao local em diferentes plataformas de borda. Como limita\c{c}\~ao, o estudo ainda depende de uma base de dados local e relativamente reduzida, o que indica a necessidade de ampliar campanhas de coleta em diferentes per\'iodos e ambientes. Como continuidade, recomenda-se expandir a valida\c{c}\~ao em outros corpos h\'idricos amaz\^onicos, incorporar novas vari\'aveis ambientais e aprofundar o uso da solu\c{c}\~ao em cen\'arios reais de monitoramento territorial.

Por fim, entende-se que o principal valor do trabalho est\'a em demonstrar uma ponte entre pesquisa aplicada, tecnologia embarcada e relev\^ancia socioambiental. Ao mostrar que o monitoramento da qualidade da \'agua pode ser realizado de forma offline, local e com potencial de uso em regi\~oes isoladas, o estudo oferece uma contribui\c{c}\~ao concreta para a discuss\~ao sobre tecnologias sustent\'aveis voltadas \`a Amaz\^onia.
